\subsection{线性映射的秩}

定义3. 设E、F是域K上的两个向量空间。对于E到F的每一个线性映射u，F的子空间u(E)的维数称为u的秩，记为rg(u)。

若 \(\mathbf{N} = \operatorname{Ker}(u)\) ，则 \(\mathbf{E} / \mathbf{N}\) 同构于 \(u(\mathbf{E})\) ，由此得关系式

\[
\operatorname{rg} (u) = \operatorname{codim} _ {\mathrm{E}} (\operatorname{Ker} (u))\tag{11}
\]

，从而

\[
\operatorname{rg} (u) + \dim (\operatorname{Ker} (u) = \dim E.\tag{12}
\]

此外，由公式(3)

\[
\operatorname{rg} (u) + \dim (\operatorname{Coker} (u)) = \dim F.\tag{13}
\]

命题9. 设E、F是域K上的两个向量空间，且u:E→F是一个线性映射。

\begin{enumerate}
\def\labelenumi{(\roman{enumi})}
\item
  \(\operatorname{rg}(u) \leqslant \inf(\dim E, \dim F)\) .
\item
  假设 \(\mathbf{E}\) 是有限维的；为了使 \(\operatorname{rg}(u) = \dim \mathbf{E}\) ，必要且充分的条件是 \(u\) 是单射的。
\item
  假设 \(\mathbf{F}\) 是有限维的；为了使 \(\operatorname{rg}(u) = \dim \mathbf{F}\) ，必要且充分的条件是 \(u\) 是满射。
\end{enumerate}

这直接由关系式(12)和(13)得出。

推论。设 E 是维数为 n 的有限维向量空间，u 是 E 的一个自同态。以下性质是等价的：

\begin{enumerate}
\def\labelenumi{(\alph{enumi})}
\item
  \(u\) 是双射；
\item
  \(u\) 是单射；
\item
  \(u\) 是满射；
\item
  \(u\) 是右可逆的；
\item
  \(u\) 是左可逆的；
\item
  \(u\) 的秩为 \(n\) .
\end{enumerate}

若 E 是无限维向量空间，则存在 E 的单射（相应地，满射）自同态，它们不是双射（习题 9）。

设 K， \(K'\) 为两个域， \(\sigma:K\to K'\) 为 K 到 \(K'\) 上的一个同构，E 为一个向量 K-空间， \(E'\) 为一个向量 \(K'\) -空间，而 \(u:K\to K'\) 为关于 \(\sigma\) 的半线性映射（§1，第13号）； \(u(\mathbf{E})\) 作为 \(E'\) 的子空间，其维数也称为 u 的秩。它也是将 u 视为 E 到 \(\sigma_{*}(\mathbf{E}')\) 的线性映射时 u 的秩，因为 \(u(\mathbf{E})\) 的每一组基也是 \(\sigma_{*}(u(\mathbf{E}))\) 的基.

